\documentclass[10pt,a4paper]{amsart}
\usepackage[margin=19mm]{geometry}
\usepackage{amsmath,amssymb,mathtools}
\usepackage[scr=rsfs,cal=euler]{mathalfa}
\usepackage{xfrac}
\usepackage[unicode,hidelinks,bookmarks=false]{hyperref}
\newcommand{\ie}{i.e.\ }
\newcommand{\cf}{cf.\ }
\newcommand{\resp}{respectively\ }
\newcommand{\Subsection}{\subsection}
\providecommand{\nobreakdash}{\nobreak\hbox{-}}
\setlength{\emergencystretch}{2em}
\multlinegap0pt
\usepackage{amssymb}
\usepackage{mathtools}
\usepackage{bm}
\usepackage[shortlabels]{enumitem}
\newtheorem{theorem}{Theorem}
\title{\textbf{Critical Structure of Axisymmetric Navier--Stokes with Swirl}}
\author{Rishad Shahmurov}
\date{Source: arXiv:2606.07869v2 (CC BY 4.0)}
\begin{document}
\maketitle

\noindent\emph{Source and licence.} \textbf{Critical Structure of Axisymmetric Navier--Stokes with Swirl}. Version arXiv:2606.07869v2 (CC BY 4.0), distributed by arXiv under the Creative Commons Attribution 4.0 International licence, http://creativecommons.org/licenses/by/4.0/, as recorded in the arXiv licence metadata for this exact version and shown on the abstract page https://arxiv.org/abs/2606.07869v2. Full licence notice: \texttt{licenses/arxiv-2606.07869-CC-BY-4.0.txt}.\par\smallskip

\noindent\textit{Editorial scope.} Counted unit: the article's theorem thm:q1213 (source0.tex lines 515--546). The article's complete original proof of it is delivered with it (source0.tex lines 547--678, one \texttt{proof} environment of 132 lines). No mathematics is written, repaired or re-derived: the statement and the proof are the article's own lines, in the article's own notation, and no retained line is substituted or re-flowed.  No further source-local context is needed: every cross-reference of every retained line resolves inside the two delivered spans.  Nothing was omitted from any retained span, so the package carries no omission record.  No retained line contains a citation, so the wrapper carries no bibliography and no bibliography entry.  No retained line contains a figure environment, an image-inclusion command or drawing code, so no image is delivered and the package declares no asset dependency.  Editorial formatting only, in addition: an AMS \texttt{amsart} preamble that restates the article's own declarations and macros used by the retained lines, namely the environments \texttt{theorem} on separate counters, together with the standard packages those lines need.  The retained environments print in this document as Theorem 1 in order of appearance, whereas the article prints the same items with the numbering of its own full document.  The complete native source of the article as distributed in the arXiv e-print for this exact version is bundled unchanged as \texttt{sources/202285/source0.tex}.
\begin{theorem}\label{thm:q1213}
On the positive-descendant source-built fresh-tangent regime with bounded
deformation, finite marks, and relative donor distance
$D_{\rm rel}=R^{q-1}$, a fixed order-one fresh tangent requires
\[
N\gtrsim D_{\rm rel}^7,
\qquad
\frac{T_{\rm prep}}{T_R}\gtrsim D_{\rm rel}^6,
\qquad
\|\widetilde G^f\|_{L^2(A_{D_{\rm rel}})}^2\gtrsim D_{\rm rel}^7.
\]
On the uncancelled visibility branch,
\[
\boxed{
\mathcal D_{G,\rm prep}
\gtrsim R D_{\rm rel}^{13}
=R^{13q-12}.}
\]
More generally,
\[
\boxed{\begin{aligned}
\text{fresh source-built preparation}\Longrightarrow{}&
\mathcal D_{G,\rm prep}\gtrsim R D_{\rm rel}^{13}\\
&\vee\ \text{same-shell cancellation}\\
&\vee\ \text{unbounded deformation}\\
&\vee\ \text{boundary/frequency loss}.
\end{aligned}}
\]
Thus $q<12/13$ is excluded only on the uncancelled, bounded-deformation,
source-built regime; $q=12/13$ is critical, while $q>12/13$ is not excluded by
this payment.
\end{theorem}

\begin{proof}
We give the preparation calculation in full because the seventh and thirteenth
powers arise from two different mechanisms: the first from the far-field
freshness detector and the second from the amount of physical time for which
that detector must remain active.

Set
\[
 \varepsilon:=D_{\rm rel}^{-1}=\frac{R}{D_{\rm phys}}.
\]
Let $W_0,W_1\ge0$ denote the two consecutive normalized positive cargo profiles,
supported in a fixed ball $B_L$, and normalize their masses so that they stay
between two fixed positive constants.  For the second Hodge derivative
\[
 K_2=\partial_{zz}\mathcal N_5,
 \qquad \mathcal N_5(x)=\frac1{8\pi^2}|x|^{-3},
\]
write
\[
 \Theta_j=K_2*W_j.
\]
Taylor expansion at distance $D_{\rm rel}\gg L$ gives
\[
 \Theta_j=Z_jK_2+O\!\left(\frac{L}{D_{\rm rel}}\right)Z_jK_2
 \quad\hbox{in }L^2(A_{D_{\rm rel}}),
\]
where $Z_j=\int W_j$.  Since $K_2$ is homogeneous of degree $-5$ in five
space dimensions,
\[
 \|K_2\|_{L^2(|x|>D_{\rm rel})}^2\asymp D_{\rm rel}^{-5}.
\]
After the leading masses are matched, the angle between the two detector
columns is therefore only of order $L/D_{\rm rel}$.  Equivalently, the small
eigenvalue of their two-column Gram matrix is of order
\[
 D_{\rm rel}^{-5}
 \left(\frac{L}{D_{\rm rel}}\right)^2
 \asymp D_{\rm rel}^{-7},
\]
with $L$ fixed.  Hence the inverse Gram matrix has norm of order at least
$D_{\rm rel}^{7}$.  A response matrix with a fixed nonzero smallest singular
value can consequently be produced only by an owner whose normalized
$L^2$-mass satisfies
\[
 \|\widetilde G^f\|_{L^2(A_{D_{\rm rel}})}^2
 \gtrsim D_{\rm rel}^{7}.
\]
This is the origin of the seventh power.

We next quantify how long a genuinely source-born owner needs in order to
acquire that response.  Let
\[
 \Gamma_\nu:=\frac{\|\Gamma\|_\infty}{\nu},
 \qquad
 \mathcal A_{\rm def}:=\int_I\|\nabla b(t)\|_\infty\,dt
\]
on the preparation interval $I$.  A donor turnover has normalized rate
\[
 m\asymp \kappa_*D_{\rm rel}^{3};
\]
the power three is the degree of the rate-leading second-Hodge response after
child normalization.  The difference of the two positive cargo detectors has
one additional far-field cancellation and hence terminal $C^1$ size
$O(\varepsilon)$.  Transporting this difference detector backwards through the
actual drift gives the source-response estimate
\[
 |\mathcal R_I|
 \le C\Gamma_\nu^2|I|\,\varepsilon
 e^{C\mathcal A_{\rm def}}.
\]
If the fresh cohort is assembled from $N$ donor turnovers, then
$|I|\lesssim N/m$.  Reaching a fixed normalized fresh response means that
$|\mathcal R_I|\gtrsim c m$.  Therefore
\[
 m
 \lesssim
 \Gamma_\nu^2\frac{N}{m}\varepsilon e^{C\mathcal A_{\rm def}},
\]
and hence
\[
 N e^{C\mathcal A_{\rm def}}
 \gtrsim
 \Gamma_\nu^{-2}m^2\varepsilon^{-1}.
\]
Using $m\asymp\kappa_*\varepsilon^{-3}$ gives
\[
 N e^{C\mathcal A_{\rm def}}
 \gtrsim c\Gamma_\nu^{-2}\varepsilon^{-7}
 =c\Gamma_\nu^{-2}D_{\rm rel}^7.
\]
On the bounded-deformation branch, $\mathcal A_{\rm def}\le A_0$, all fixed
constants may be absorbed, yielding
\[
 N\gtrsim D_{\rm rel}^7.
\]
One donor turnover occupies a fraction $D_{\rm rel}^{-1}$ of a child diffusion
clock.  Consequently
\[
 \frac{T_{\rm prep}}{T_R}
 \gtrsim ND_{\rm rel}^{-1}
 \gtrsim D_{\rm rel}^{6}.
\]
This is the origin of the sixth power in the preparation age.

It remains to convert normalized current mass and normalized duration back to
physical dissipation.  In the chosen child normalization, one child diffusion
clock carrying normalized $G$-mass $M_G$ contributes physical
azimuthal-vorticity dissipation of order $R M_G$.  Therefore the last ramp
costs at least
\[
 \mathcal D_{G,\rm prep}
 \gtrsim
 R\,(D_{\rm rel}^7)(D_{\rm rel}^6)
 =R D_{\rm rel}^{13}.
\]
Since $D_{\rm rel}=R^{q-1}$,
\[
 R D_{\rm rel}^{13}
 =R^{1+13(q-1)}
 =R^{13q-12}.
\]
Thus the cost diverges for $q<12/13$, is scale invariant for $q=12/13$, and
can be Zeno-summable for $q>12/13$.

Finally, this physical lower bound applies only if the total current retains a
fixed fraction of the freshly generated current.  If it does not, the triangle
inequality forces an older or incoming component of comparable size on the
same donor shell.  That is the stated same-shell cancellation alternative.
Loss of the bounded-deformation or fixed localization hypotheses gives the
remaining deformation, boundary, contact, or frequency alternatives.  No such
failure is counted as dissipation payment.  This proves the theorem.
\end{proof}

\end{document}
